% capitulos/cap5/operacao.tex -- §5.13 em diante: implantação, desempenho e síntese do capítulo.
% Parte do Capítulo 5 (capitulos/cap5_arquitetura.tex); os rótulos \label e os números de seção não mudaram.

\section{Implantação e Operação}\label{sec:arq-implantacao}

Por decisão do Engenheiro Chefe (\decisao{D-04}, \adr{ADR-011}), o alvo de implantação é uma VPS tradicional com Docker Engine e Docker Compose. A aplicação é empacotada em uma imagem construída em dois estágios (instalação e \en{build}; execução) que executa o servidor Next.js (\texttt{next start}) como usuário sem privilégios, em um único contêiner. O arquivo do SQLite fica em um \en{bind mount} do hospedeiro (\arq{./data}, montado em \arq{/data}), fora da imagem, e as variáveis de ambiente vêm de um arquivo \arq{.env} também fora da imagem. Um \en{proxy} reverso com TLS (Caddy ou Nginx) encaminha as requisições para a porta 3000 do contêiner, publicada apenas em \arq{127.0.0.1} (Figura~\ref{fig:impl-implantacao}). O passo a passo, com o \arq{Dockerfile}, o \arq{docker-compose.yml} e o \arq{.dockerignore} do repositório, está no Apêndice~\ref{ap:implantacao}.

\vitrinefigura[tipo=diagrama, arquivo=impl-implantacao, agente=Arquiteto, largura=0.95,
  origem={doc/diagrams/src/impl-implantacao.puml},
  descricao={Diagrama de implantação: VPS com Docker Engine; proxy reverso com TLS; contêiner vitrine com o build do Next.js, o cliente do Prisma e o adaptador better-sqlite3; volume ./data com o arquivo SQLite (e WAL) e cópias de segurança; arquivo .env fora da imagem.}]
  {Diagrama de implantação}{fig:impl-implantacao}

Os requisitos normativos sobre a imagem, atendidos pelo módulo M13 (merge M-008) e válidos para qualquer mudança nela, são:

\begin{itemize}
  \item \textbf{Módulo nativo e OpenSSL.} O \arq{better-sqlite3} é um módulo nativo (\adr{ADR-014}); os estágios de \en{build} e de execução usam a mesma imagem-base, na mesma versão do Node.js do desenvolvimento (hoje, \texttt{node:24-bookworm-slim}, com \en{glibc}), e o estágio de \en{build} precisa de compilador (\texttt{python3}, \texttt{make}, \texttt{g++}) caso não haja binário pré-compilado. A imagem \texttt{slim} \emph{não} traz a biblioteca OpenSSL de que o motor de migrações do Prisma depende; por isso o \texttt{openssl} é instalado nos dois estágios (com ele, o \texttt{prisma migrate deploy} funcionou no contêiner; a imagem sem ele não foi testada);
  \item \textbf{Cliente gerado.} \texttt{prisma generate} roda antes de \texttt{next build}, pois \arq{src/generated/prisma} não é versionado;
  \item \textbf{Migrações como passo explícito.} \texttt{prisma migrate deploy} é executado por um comando próprio antes de subir a nova versão (por exemplo, \texttt{docker compose run -{}-rm vitrine npx prisma migrate deploy}), e não pelo comando de partida da aplicação; isso exige que a CLI do Prisma e o \arq{dotenv} estejam na imagem, o que está garantido desde o merge M-002, que os moveu para \texttt{dependencies} (\adr{ADR-014});
  \item \textbf{Administrador inicial.} Depois das migrações e do \en{seed}, o administrador é criado por \texttt{docker compose run -{}-rm vitrine npm run admin:create}, que exige \arq{ADMIN_EMAIL} e \arq{ADMIN_PASSWORD} (e aceita \arq{ADMIN_NAME}) no ambiente da execução. O \en{seed} (\texttt{prisma db seed}) \emph{não} cria o administrador: ele só garante a configuração padrão (\adr{ADR-018});
  \item \textbf{\texttt{tsx} na imagem.} O \texttt{admin:create} e o \texttt{prisma db seed} rodam com o \texttt{tsx}, hoje em \texttt{devDependencies}; a imagem de produção deve manter as dependências de desenvolvimento (\texttt{npm ci} completo) ou mover o \texttt{tsx} para \texttt{dependencies}, caso contrário esses dois comandos deixam de funcionar;
  \item \textbf{\texttt{DATABASE\_URL} no \en{build}.} O \texttt{next build} importa as rotas, que importam \arq{db.js}, e este lança erro se \arq{DATABASE_URL} não está definida; o estágio de \en{build} do Dockerfile deve defini-la com qualquer valor bem formado (por exemplo, \texttt{ENV} \arq{DATABASE_URL=file:/data/vitrine.db}), pois nenhuma consulta ocorre no \en{build} e o arquivo não precisa existir --- nem o diretório. Isso é \emph{garantido por código}, e não por acaso: \texttt{getSiteConfiguration} (\arq{src/app/_lib/site.js}) devolve \texttt{null} sem consultar quando \arq{NEXT_PHASE} vale \texttt{phase-production-build}, e todo \en{layout} ou página que possa ser pré-renderizado deve ler dados do banco por essa guarda ou por outra equivalente; sem ela, o \texttt{next build} falha ao pré-renderizar \arq{/_not-found} (\adr{ADR-021}; Seção~\ref{sec:arq-tema-estatico}). As páginas estáticas saem do \en{build} com os padrões do CSS e passam à configuração do banco na primeira revalidação, o que se repete a cada recriação do contêiner (comportamento por projeto, \en{by design}; Seção~\ref{sec:arq-tema-estatico}). Segredos (\arq{BETTER_AUTH_SECRET}, \arq{ADMIN_PASSWORD}) \emph{não} entram no \en{build}: o aviso do segredo padrão registrado nessa etapa é esperado e inofensivo;
  \item \textbf{Segredo em execução.} \arq{BETTER_AUTH_SECRET} e \arq{BETTER_AUTH_URL} (em \texttt{https}, para o \en{cookie} \texttt{Secure}) são obrigatórios na execução; sem o segredo, as rotas de autenticação e de API respondem 500;
  \item \textbf{Instância única.} Um único contêiner, pois o cache de ISR é por instância (Seção~\ref{sec:arq-isr}), o SQLite tem um escritor por vez e a limitação de tentativas de login guarda os contadores em memória;
  \item \textbf{Usuário sem privilégios e arquivos graváveis.} O contêiner executa como \texttt{node} (UID 1000), nunca como \texttt{root}. A aplicação copiada para a imagem (\texttt{COPY -{}-chown}) e o diretório \arq{/data} pertencem a esse usuário; a ISR regenera páginas em \arq{.next/server/app/} (e usa \arq{.next/cache}), de modo que todo o \arq{.next} precisa ser gravável. No hospedeiro, o diretório \arq{data/} deve pertencer ao UID 1000;
  \item \textbf{Partida sem intermediário.} O comando é \texttt{node\_modules/.bin/next start}, sem o npm no meio, e o \texttt{init: true} do \en{compose} (o \texttt{docker-init} como processo 1) repassa os sinais de parada e recolhe processos órfãos;
  \item \textbf{Rede.} A porta 3000 é publicada somente em \arq{127.0.0.1}; o acesso público é do \en{proxy} reverso;
  \item \textbf{Cabeçalhos de segurança.} A CSP e os demais cabeçalhos (Seção~\ref{sec:arq-csp}) vêm da própria aplicação e não dependem do \en{proxy}, que não deve alterá-los;
  \item \textbf{Espaço em disco.} A construção da imagem exige pelo menos 6~GB livres no disco do Docker (Apêndice~\ref{ap:implantacao}, Seção ``Pré-requisitos'');
  \item \textbf{Segredos.} \arq{.env} e \arq{data/} constam do \arq{.dockerignore} e do \arq{.gitignore};
  \item \textbf{Escopo multiusuário (M15 a M17; Apêndice~\ref{ap:implantacao}).} O banco roda em WAL, com os arquivos \texttt{-wal} e \texttt{-shm} ao lado dele, em volume \emph{local} (nunca de rede), e a cópia de segurança continua sendo o \texttt{VACUUM INTO} (Seção~\ref{sec:arq-wal}); as variáveis \texttt{MAX\_CONTENTS\_PER\_USER}, \texttt{MAX\_BODY\_BYTES} e, opcionalmente, \texttt{ADMIN\_HANDLE} entram no \arq{.env}; o \en{proxy} reverso aplica limite de taxa por endereço às rotas públicas, como mitigação do documento inexistente em espaço válido (\adr{ADR-030}), e o disco do volume do \arq{.next} é monitorado; o \texttt{admin:reset-password} recupera uma conta e o \texttt{admin:create} recupera o acesso sem administrador ativo (RN22).
\end{itemize}

O relatório de IC citava a plataforma Vercel como exemplo de implantação \cite{reis2026ic}. O sistema de arquivos de funções \emph{serverless} é efêmero, de modo que um arquivo SQLite local não persistiria entre invocações nesse tipo de ambiente; a decisão por uma VPS com volume elimina esse obstáculo, que constava como risco R01, para esta versão. Plataformas dessa natureza só seriam viáveis com a troca do \en{driver adapter} por um banco remoto compatível (por exemplo, libSQL/Turso), o que está fora do escopo (\adr{ADR-011}).

\section{Desempenho e Escalabilidade}\label{sec:arq-desempenho}

O perfil de carga é dominado por leituras. A ISR evita consultas ao banco na maioria das requisições públicas; quando há regeneração, a consulta da árvore é linear no número de nós e a resolução do caminho é proporcional à profundidade. O SQLite em modo WAL permite leitores simultâneos à escrita, que é serializada \cite{sqlitewal}; como as escritas ocorrem por ação de poucos autores (de um a algumas dezenas, cada um limitado a 500 conteúdos e a um limite de taxa; RN25), a limitação de um escritor por vez não é relevante para o perfil esperado \cite{sqlitewhentouse}, e a latência de cada escrita é a do \texttt{fsync} do disco (Seção~\ref{sec:arq-wal}). A invalidação por espaço (\adr{ADR-030}) tem custo desprezível: cerca de 500 chamadas por evento, abaixo de 10~ms. Os limites práticos --- tamanho da árvore de um dono, número de documentos pré-renderizados, o crescimento do \arq{.next} com páginas inexistentes e o custo de o Proxy consultar o documento --- serão medidos no Capítulo~\ref{cap:testes} e pelo Core, na implementação do M17.

\section{Síntese do Capítulo}

O capítulo especificou o monólito modular, suas visões, regras de dependência, modelo de dados (com o estado de publicação e os links sociais), autenticação em camadas, estratégia de cache e de visibilidade, pipeline de conteúdo, tematização, contrato de API, segurança e implantação por imagem Docker. Com a ampliação de escopo (\adr{ADR-023} a \adr{ADR-031}), especificou também, como \emph{alvo} a implementar pelos módulos M15 a M20, o modelo de dados multiusuário, a autorização por dono, o espaço público por usuário com revalidação por caminhos literais, a aparência por proprietário com paletas e neutros por modo e as notificações; o texto distingue em cada seção o que está integrado do que é alvo. As observações sobre o esquema atual (Tabela~\ref{tab:obs-schema}) e os valores padrão da \adr{ADR-017} formam a lista de tarefas de refinamento dos agentes; o Capítulo~\ref{cap:implementacao} acompanha a execução.
